\chapter{Tableaux de bord et alertes}

\section{Principes de présentation}

Un seul tableau de bord pour tous les utilisateurs aurait accumulé les
indicateurs jusqu'à devenir illisible. Trois vues ont donc été créées. Leur
vocabulaire est en français et leurs requêtes renvoient un état explicite
lorsqu'aucune donnée n'est encore disponible.

\section{Vue opérateur}

La vue opérateur doit être comprise rapidement, même à distance. Elle contient
l'état de l'imprimante, le dernier défaut, les niveaux d'encre et de solvant,
le délai avant action et le nombre de pare-brise marqués par période de cinq
minutes.

\begin{table}[H]
\centering
\caption{Correspondance entre état technique et texte affiché}
\begin{tabularx}{0.85\textwidth}{p{0.33\textwidth}Y}
\toprule
Classification & Libellé opérateur \\
\midrule
\texttt{producing} & EN PRODUCTION \\
\texttt{downtime} & À L'ARRÊT \\
\texttt{idle\_jet\_on} & EN ATTENTE \\
\texttt{idle\_no\_order} & SANS COMMANDE \\
absence de mesure & AUCUNE DONNÉE \\
\bottomrule
\end{tabularx}
\end{table}

Le dernier défaut est traduit en français dans la requête. S'il n'existe
aucun défaut, le panneau affiche ``Aucun défaut'' au lieu de ``No data''.

\begin{figure}[H]
  \centering
  \includegraphics[width=\textwidth]{grafana-operateur.png}
  \caption{Tableau de bord opérateur alimenté par la simulation}
  \label{fig:grafana-operateur}
\end{figure}

\section{Vue maintenance}

La vue maintenance rassemble les informations nécessaires à un diagnostic :

\begin{itemize}
  \item codes natifs du jet et de l'impression, vitesse moteur, pression et
        viscosité générés par le simulateur 9450c ;
  \item capacités de 800~mL et 1,075~L, réserve interne de 24~heures et
        températures électronique, encre et tête ;
  \item tableau des consommables avec lot, niveau, jours avant épuisement,
        date effective de péremption et facteur limitant ;
  \item courbes d'encre et de solvant ;
  \item viscosité comparée à la consigne ;
  \item pression et température de la tête ;
  \item heures restantes avant service ;
  \item liste des défauts récents ;
  \item Pareto du temps perdu par code de défaut ;
  \item historique des lots installés.
\end{itemize}

Le tableau des consommables met côte à côte les deux échéances. Une cartouche
peut être encore remplie à 60~\% et devoir être remplacée cinq jours plus tard
à cause de sa date après ouverture. Afficher uniquement le niveau masquerait ce
cas.

\begin{figure}[H]
  \centering
  \includegraphics[width=\textwidth]{grafana-simulation-9450c.png}
  \caption{Paramètres natifs produits en direct par la simulation 9450c}
  \label{fig:grafana-simulation-9450c}
\end{figure}

\begin{figure}[p]
  \centering
  \includegraphics[width=\textwidth,height=0.78\textheight,keepaspectratio]{grafana-maintenance.png}
  \caption{Tableau de bord de maintenance alimenté par la simulation 9450c}
  \label{fig:grafana-maintenance}
\end{figure}

\section{Vue direction}

La vue direction évite les mesures brutes. Elle présente la disponibilité du
jour, le nombre de pare-brise marqués, le rendement au premier passage, le
temps d'arrêt, la répartition du temps et les principales causes d'arrêt. Le
graphique de répartition conserve une catégorie ``sans commande'', car cette
information explique la différence entre temps calendaire et temps demandé.

Si aucune demande n'a encore été observée dans la journée, le panneau de
disponibilité affiche ``NON CALCULÉE''. Une valeur de 0~\% aurait une autre
signification : la ligne aurait demandé des impressions sans aucune
production.

\begin{figure}[H]
  \centering
  \includegraphics[width=\textwidth]{grafana-direction.png}
  \caption{Tableau de bord de direction alimenté par les données simulées}
  \label{fig:grafana-direction}
\end{figure}

\section{Règles d'alerte}

\begin{longtable}{p{0.29\textwidth}p{0.45\textwidth}p{0.18\textwidth}}
\caption{Règles provisionnées dans Grafana}\\
\toprule
Alerte & Condition & Destinataire \\
\midrule
\endfirsthead
\toprule
Alerte & Condition & Destinataire \\
\midrule
\endhead
Consommable bientôt vide &
Temps estimé inférieur à 8 heures pendant 5 minutes &
Maintenance \\
Niveau bas &
Niveau inférieur à 20~\% pendant 5 minutes &
Maintenance \\
Consommable proche de la péremption &
Moins de 14 jours avant la date effective &
Maintenance \\
Consommable périmé &
Date effective dépassée &
Maintenance et qualité \\
Arrêt avec demande &
Classification \texttt{downtime} pendant 2 minutes &
Opérateur \\
Viscosité hors plage &
Écart supérieur à 0,3 pendant 15 minutes &
Maintenance \\
Absence de données &
Dernière mesure reçue depuis plus de 120 secondes &
Ingénierie \\
\bottomrule
\end{longtable}

Les alertes liées à l'arrêt utilisent le signal de demande. En revanche,
l'absence totale de télémétrie reste un incident, même sans commande : elle
peut indiquer une perte de communication ou un service arrêté.

\section{Lisibilité et langue}

Les titres et les états visibles ont été traduits en français. Les noms de
mesures, champs et codes de défaut restent techniques, car ils servent de
contrat avec l'ingestion. La traduction intervient au niveau de la requête ou
de l'affichage, sans modifier la donnée brute.

Les figures~\ref{fig:grafana-operateur} à~\ref{fig:grafana-direction}
présentent les vues de l'instance Grafana de démonstration alimentée par les
données InfluxDB. Les panneaux disposent d'un état de repli explicite afin de
ne pas confondre une absence de défaut avec une absence de données.
